\runninghead{The Lamb Receives the Kingdom}
\properlane{lamb-receives-the-kingdom}{The Slain Lamb Receives the Kingdom That Was Always His}

The kingship this Mass acclaims belongs to the Son by nature, and it is also
received by him: received in the humanity he assumed, asked for and won by his
death, and given at his exaltation. What he receives is the nations he has
redeemed, and the place where that received kingship is exercised is the
sacrifice the Mass offers. The reading begins from the Introit's verb and the
Offertory's command, and it lets the whole formulary take its order from the
Preface, where the anointed Son first offers himself on the cross and only
then hands over a kingdom.

\subsection{What the Lamb receives}
\label{sec:lamb-question}

The angels of the Introit are not praising God in the abstract. They stand
round the throne at the moment the Lamb has taken the sealed book that no one
else was worthy to open, and the reason the elders have just given is a
history: ``thou wast slain and hast redeemed us to God, in thy blood''
(Apoc 5:9). Their acclamation names seven things the Lamb is worthy to
receive, and the first two that the Introit sings are \latin{virtutem, et
divinitatem}: power, and divinity. The Offertory raises the same difficulty
from another book. The king set on Sion recites the decree, ``Thou art my son,
this day have I begotten thee,'' and the Lord goes on, \latin{Postula a me, et
dabo tibi gentes hereditatem tuam}: ``Ask of me, and I will give thee the
Gentiles for thy inheritance'' (Ps 2:7--8). If the Lamb is the eternal Son,
how can he receive divinity? If the one addressed is the Son begotten, why
must he ask for anything? The Alleluia sings only the possessing half of Dan
7:14, ``his power is an everlasting power,'' but the verse it comes from begins
``And he gave him power.''

The question was not idle in the Church's history. St Jerome, expounding that
verse of Daniel, writes that the Son of man's receiving power, honour and
kingdom must be understood by the Apostle's ``Who being in the form of God
\dots\ emptied himself, taking the form of a servant \dots\ becoming obedient
unto death'' (Phil 2:6--8), which the Arians, he says, should have read
piously (\work{In Dan.}\ II, on 7:13--14). Read without that key, texts of
receiving make the Son a creature promoted for his merits. Read with it, they
speak of what the eternal Son receives in the nature he took. The Fathers
answer the question where the second psalm and Daniel raise it, and their
answers carry over to the Lamb of the Introit, who is worthy to receive
because he was slain.

\subsection{St Hilary: the risen Son asks back the glory he had}
\label{sec:lamb-hilary}

St Hilary of Poitiers reads the decree of Psalm~2 on apostolic authority. St
Paul at Antioch had said that God ``hath fulfilled \dots\ raising up Jesus, as
in the second psalm also is written: Thou art my Son: this day have I begotten
thee'' (Acts 13:33). So Hilary takes the ``today'' of the decree as the
begetting of the firstborn from the dead (\work{Tractatus super Psalmos}
2.30). The risen Lord's own words, ``All power is given to me in heaven and in
earth'' (Mt 28:18), show that he asked for what he received; and when
the psalm goes on, \latin{Posce a me}, Hilary concludes, \latin{accepit ergo
hereditatem gentium, quam poposcit}, he received the inheritance of the
nations, which he asked for (2.31). The asking itself Hilary finds in the prayer of the Supper, ``Father, the hour is come. Glorify
thy Son \dots\ As thou hast given him power over all flesh'' (Jn 17:1--2).

What was asked, and what was given? Hilary answers with the hymn of
Philippians. He who abode in the form of God took the form of a servant, and
after the death of the cross he is exalted into the name that is above every
name: \latin{in deum namque, quia nullum ultra deum nomen est, prouehitur,
eique hoc petenti, id est, ut esset, quod fuerat ante, donatur}, for he is
raised into God, since there is no name beyond God, and to him who asks this,
namely that he might be what he had been before, it is given (2.33). The
receiving is the glorifying of the humanity the Son assumed. The Son loses nothing and gains
nothing as God; the servant he became is raised into the glory the Word always
had. And the inheritance is not territory. It is the nations baptized and
taught, carried over \latin{ex iniusto \dots\ iure dominantium in regnum
aeternum diuinumque}, out of the unjust right of those who lorded it over them
into an eternal and divine kingdom, the inheritance \latin{primogeniti huius ex
mortuis}, of this firstborn from the dead (2.31). Hilary is expounding the
psalm, not the Epistle, but his words for the inheritance are the Epistle's:
those transferred ``into the kingdom of the Son of his love,'' and the Son
``the firstborn from the dead'' (Col 1:13, 18). Later in the same tractate he
draws the consequence for the inheritance itself: the kingdom of God is where
sin is overcome and death destroyed, and those who become co-heirs of the
eternal King must themselves become kings (2.42--43).

\subsection{St Augustine and St Thomas: one dominion by two rights}
\label{sec:lamb-two-rights}

St Augustine divides the two verses differently. ``To-day have I begotten
Thee,'' though it may seem spoken of the birth in the flesh, he gives ``a
divine interpretation'': it is the eternal generation, and ``today'' signifies
the presence of eternity (\work{Enarrationes in Psalmos} 2.6). Only verse~8 is
temporal: ``Ask of Me'' ``has at once a temporal sense with reference to the
Manhood which He took on Himself, who offered up Himself as a Sacrifice in the
stead of all sacrifices, who also maketh intercession for us'' (2.7). So
Augustine and Hilary agree on verse~8, the asking and receiving of the nations
by the Son made man, and differ on verse~7: Hilary reads the decree of the
resurrection, Augustine of the eternal generation.

St Thomas Aquinas states Augustine's division as a principle. The two verses
give Christ dominion over the nations \latin{duplici jure: scilicet
hereditario, et meritorio}, by a double right, of inheritance and of merit.
The right of inheritance is the eternal generation, \latin{hodie genui te,
idest aeternaliter}. The right of merit is the asking: God wills that we
receive his gifts by asking, \latin{et hoc exemplum voluit ostendere per
Christum, quia voluit quod peteret, quod sibi jure hereditario competebat}, and
he willed to show this example in Christ, willing that he should ask for what
belonged to him by right of inheritance. Christ asks by his prayer, ``for them
also who through their word shall believe in me'' (Jn 17:20), and by his
passion, the death that redeems those called to the eternal inheritance (Heb
9:15); and he holds the nations not as ministers hold them, as Peter or Paul
did, but as lord, the Son in his own house (Heb 3:5--6; \work{Postilla super
Psalmos} 2, nn.~5--6). The double right
removes the scandal of the Offertory. The Son does not ask because he lacks;
he asks because the Father wills that even what is owed be received as a gift,
and because the asking is his passion.

Cassiodorus gives the rule its widest form. Whenever the Lord Christ asks that
something be given him, or another asks it for him, \latin{naturam humanitatis
evidenter ostendi}, the nature of his humanity is plainly shown; for the
almighty Word, having all that is the Father's, needs nothing to be given
(\work{Expositio Psalmorum} 71, on verse~1, with Jn 16:15). On Ps 2:8 he applies
it: \latin{Quidquid enim in tempore accepit Christus, secundum hominem ea quae
non habebat acquirit}; the possession is given because \latin{natura
humanitatis accepit, quod semper divina possedit}, the human nature received
what the divine always possessed. Theodoret of Cyrus, on the same verse, says
that he who is Lord by nature as God \latin{etiam ut homo universorum dominium
accipit}, receives also as man the lordship of all; on the Introit's verse he
adds that the psalm teaches both, that as God he has authority over all and as
man receives it from his own Father (\work{Interpretatio in Psalmos} 2; 71).
Theodoret parts from Augustine on verse~7, which he will not apply to the
Godhead at all, since of that generation the Father says ``from the womb before
the day star I begot thee'' (Ps 109:3): as man Christ receives that voice too.

St Robert Bellarmine hears three generations in verse~7, the eternal, the
birth in the flesh and the resurrection, and makes the first the ground of the
decree: ``to Christ, as being the true and natural Son of God, is due all
power in heaven and on earth.'' On verse~8 the Father says to him, ``You my
natural Son, the incarnation of my power raised from the dead, have just right
to ask me for power over all nations''; the kingdom so given is ``his spiritual
kingdom, that is, his Church,'' not bounded like David's or Solomon's but over
all nations, ``infidels even included'' (\work{Explanatio}, Ps~2).

\subsection{St Bede: the Lamb and the Lion, King and Priest}
\label{sec:lamb-bede}

St Bede reads the Apocalypse itself in the same grammar. The one whom the
ancient calls a Lion and John sees as a Lamb is one Lord: a Lamb in dying
innocently, a Lion in conquering death. When the Lamb takes the book from the
right hand of him who sits on the throne (Apoc 5:7), \latin{dispensationem
ipsam Incarnationis a Patre et a seipso, secundum id quod Deus est, dispositam
accepisse Filius hominis dicitur}: the Son of man is said to receive the very
dispensation of the Incarnation, which the Father and he himself, according as
he is God, had disposed. He is the Lamb in his humanity and the Father's right
hand in his deity (\work{Explanatio Apocalypsis} I, on 5:5--7). The one who
gives and the one who receives are the same Son, in his two natures. On the
greeting of chapter~1, Bede reads ``the first begotten of the dead'' as the one
who for the suffering of death was crowned with glory and honour (Heb 2:9; Phil
2:9), and on the clause of Apoc 1:6 that the Introit leaves out, ``hath made us
a kingdom, and priests,'' he gives the consequence for the redeemed: because
the King of kings and heavenly Priest united us to his body by offering
himself for us, no saint lacks a share in the priesthood
(see the Preface, above).

Victorinus of Pettau had read the sealed book in the same direction. To open it
is to overcome death for man; because the Lamb overcame death as testator, it
was just that he be appointed heir and possess the human members (on Apoc
5:2--5). The Lamb receives an inheritance by dying, and the inheritance is
mankind.

\subsection{The objection: he received nothing in addition}
\label{sec:lamb-objection}

The strongest objection comes from St John Chrysostom, and it comes from this
Mass's own Gospel. On ``For this was I born,'' Chrysostom writes: ``If then He
was born a king, all His other attributes are by Generation, and He has
nothing which He received in addition.'' He draws at once a rule for every text
that speaks of giving to the Son: ``So that when you hear that, As the Father
has life in Himself, so has He given to the Son also to have life, deem of
nothing else but His generation, and so of the rest'' (\work{Homilies on John}
83.4; Jn 5:26). St Cyril of Alexandria says of the same verse that Christ's
sovereignty is ``not \dots\ added unto it from without, but \dots\ essentially and
naturally inherent'' (\work{Commentary on John} XII). A reading that made
Christ king only from Easter, or that let the Introit's \latin{accipere
divinitatem} mean that the Lamb becomes divine, would contradict them and the
faith they defend.

The answer that Augustine's division gives grants Chrysostom's premise. The
kingship by generation is the decree of verse~7, and it is not received; the
dominion asked for in verse~8 is the same dominion held now in the humanity
and won by the passion. But the answer does not meet Chrysostom on his own
ground. His rule refers a text of giving to the eternal generation; the rule
Cassiodorus states, and Theodoret applies, refers the asking and receiving to
the humanity. Chrysostom does not expound Ps 2:8 at this place, and his rule
is not a reading of it. The witnesses agree that Christ is King by nature, and
they read the texts of giving and asking addressed to the Son by two different
rules, and on that narrower question they disagree. Hilary's reading of verse~7 of the
resurrection is compatible with the answer only because he does not make the
resurrection the beginning of the Sonship: what is given is \latin{ut esset,
quod fuerat ante}. Bellarmine's first generation is the eternal one.

The instituting act of the feast holds both titles together. \work{Quas
primas} quotes St Cyril for a dominion that is Christ's by essence and by
nature, grounded in the hypostatic union, so that to him as man angels and men
are subject, and then names a second title, acquired by redemption, citing
1~Pet 1:18--19 on the price of his precious blood (n.~13).

\subsection{How the whole Mass receives the kingdom}
\label{sec:lamb-elements}

Read in this way, every element of the formulary takes a place in one
movement, from the eternal right to the received dominion and its exercise at
the altar.

The \textbf{Introit} is the reading's text: the Lamb \latin{qui occisus est}
is worthy to receive because he was slain, and \latin{Ipsi gloria et imperium}
is what he receives. Its \latin{accipere \dots\ divinitatem}, read with
Hilary's servant raised into the name of God and Bede's Son of man who
receives what he disposed as God, means on this reading that the slain Lamb
receives, in the nature he took, the name and glory of the God he always was.
He does not become what he was not; as man he is given to be what he had been
before. Schuster's note on the antiphon, that the exaltation is
in proportion to the humiliation and that John names the Passion first, states
the same order for the Introit alone (p.~201). Its verse, \latin{Deus,
iudicium tuum Regi da}, is itself a giving to the King and the King's Son,
which Augustine reads of the Father giving all judgment to the Son (Jn 5:22)
and Theodoret of the authority he has as God and receives as man. The
\textbf{Collect} confesses that the Father willed, \latin{voluisti}, to
restore all things in his Son, and asks that the nations be subjected to him:
the inheritance asked for in the psalm, still being given. The
\textbf{Epistle} supplies the means and the moment of the receiving: redemption
``through his blood,'' peace ``through the blood of his cross,'' the
``firstborn from the dead'' who is head of the Church according to the
humanity in Theodoret's reading and \latin{secundum quod huiusmodi} in
Aquinas's.

The \textbf{Gradual} sings the inheritance in possession, the dominion from sea
to sea and the homage of the kings. The \textbf{Alleluia} keeps only the power
as possessed, \latin{potestas aeterna}; the giving that Jerome expounds by
Philippians stands in the half of the verse that is not sung. In the
\textbf{Gospel} the King stands bound before the governor, the passion already
begun. Chrysostom says he was made captive ``by consenting, and giving Himself
up'' (\work{Homilies on John} 83.4): the kingship of verse~8 is being acquired
in the very act of the trial, by the prisoner who could have had servants
fight for him. The \textbf{Offertory} is the reading's centre, the command to
ask, sung as the gifts are brought to the altar; Schuster, on this antiphon,
calls it the deed of investiture received from the Father (p.~203).

The \textbf{Secret} then names the asking that the Fathers found in verse~8.
\latin{Hostiam \dots\ humanae reconciliationis offerimus}: the sacrifice by
which, in Augustine's words, he ``offered up Himself as a Sacrifice in the
stead of all sacrifices'' is offered now, \latin{quem sacrificiis praesentibus
immolamus}, and what is asked of him is the gift that Hilary's inheritance
consists of, the nations' unity and peace. The \textbf{Preface} gives the
reading its order. The Son is anointed \latin{Sacerdotem aeternum et
universorum Regem}; he offers himself \latin{in ara Crucis} so as to
accomplish the mysteries of man's redemption; only then, \latin{suo subiectis
imperio omnibus creaturis}, does he hand over the eternal and universal
kingdom. The \textbf{Communion} sings the session of the exalted King,
\latin{Sedebit Dominus Rex in aeternum}. And the \textbf{Postcommunion} asks
for those who fight under his standards a share in his reign: Hilary's
co-heirs who must become kings, Bede's members who share the priesthood of the
King who offered himself. The standard they fight under is, in Schuster's note
on that prayer, the Cross (p.~204).

\begin{fourSenses}
\item[Literal.] An apocalyptic hymn acclaims the slain Lamb as worthy of power;
a royal psalm prays that God give judgment to the king and the king's son,
which Bellarmine reads literally of David's prayer for Solomon and Augustine,
Cassiodorus and Theodoret of Christ alone; another psalm has God promise the
nations to his Son, who is to ask for them; Paul thanks the Father for
redemption through the blood of the Son's cross; Jesus, bound, acknowledges
before the governor that he is a king.
\item[Allegorical.] Christ, God from eternity and King by that right, receives
as man, through his passion and resurrection, dominion over the nations, and
the nations so received are the Church gathered from them. The Lamb who dies
and the Lion who conquers are one Lord, who as God disposes the economy he
receives as man (Bede).
\item[Moral.] Kingship is received through obedience unto death, and God wills
that even what is owed be asked (Aquinas). The faithful come into the kingdom
in the same way: by asking, and by offering themselves with the Priest-King as
members who share his priesthood (Bede); and they reign as co-heirs by
overcoming sin in themselves (Hilary).
\item[Anagogical.] The firstborn from the dead is the first of the exalted.
His members will be raised after him and will reign with him for ever in the
heavenly seat for which the Postcommunion asks.
\end{fourSenses}
